\section{Results and Analysis}

\subsection{Quantitative Results}

The LiDAR-only PointPillars reference and the attention-bottleneck fusion
detector were evaluated on the 3,769-sample KITTI validation split using
AP$_{40}$. Both models used the same validation split, object classes,
point-cloud range, effective batch size, random seed, and KITTI evaluator.
The best checkpoint for each model was selected using overall 3D AP$_{40}$ at
moderate difficulty. The LiDAR-only reference reached its best result at epoch
28, while the fusion detector reached its best result at epoch 20.

\begin{table}[ht]
\centering
\caption{Overall KITTI validation AP$_{40}$.}
\label{tab:overall_results}
\begin{tabular}{lrrrrrr}
\hline
& \multicolumn{3}{c}{BEV AP$_{40}$} &
\multicolumn{3}{c}{3D AP$_{40}$} \\
Model & Easy & Moderate & Hard & Easy & Moderate & Hard \\
\hline
LiDAR-only PointPillars
& 64.5245 & 52.8874 & 49.6791 & 49.5614 & 38.6868 & 36.1066 \\
Attention-bottleneck fusion
& \textbf{66.3382} & \textbf{54.6478} & \textbf{51.9588}
& \textbf{54.4341} & \textbf{42.6696} & \textbf{39.9782} \\
\hline
Difference
& +1.8137 & +1.7604 & +2.2797 & +4.8727 & +3.9828 & +3.8716 \\
\hline
\end{tabular}
\end{table}

The attention-bottleneck fusion detector outperforms the LiDAR-only reference
at every BEV and 3D difficulty level. At moderate difficulty, it improves BEV
AP$_{40}$ by 1.7604 points and 3D AP$_{40}$ by 3.9828 points. The latter is a
10.29\% relative improvement over the LiDAR-only result.

\begin{table}[ht]
\centering
\caption{Per-class 3D AP$_{40}$ at moderate difficulty.}
\label{tab:per_class_results}
\begin{tabular}{lrrr}
\hline
Class & LiDAR-only & Fusion & Difference \\
\hline
Pedestrian & 30.3730 & \textbf{32.2419} & +1.8689 \\
Cyclist & 34.1321 & \textbf{37.7407} & +3.6086 \\
Car & 51.5553 & \textbf{58.0263} & +6.4710 \\
\hline
\end{tabular}
\end{table}

Fusion improves moderate 3D AP$_{40}$ for all three classes. The largest
absolute improvement is observed for Car, followed by Cyclist and Pedestrian.
The best fusion checkpoint obtains a moderate 3D AP$_{40}$ of 42.6696 at epoch
20, compared with 41.5993 at epoch 40. This decrease indicates mild
overfitting during the second half of training and supports validation-based
checkpoint selection.

\subsection{Camera Ablation}

The best fusion checkpoint was reevaluated with correct camera images,
all-black images, and a deterministic seed-0 derangement of the validation
images. In the shuffled condition, every LiDAR sample receives an image from
a different validation sample while its LiDAR points and annotations remain
unchanged.

\begin{table}[ht]
\centering
\caption{Camera-input ablation using the epoch-20 fusion checkpoint.}
\label{tab:camera_ablation}
\begin{tabular}{lrr}
\hline
Camera input & Moderate BEV AP$_{40}$ & Moderate 3D AP$_{40}$ \\
\hline
Correct image & \textbf{54.6478} & \textbf{42.6696} \\
Black image & 54.1208 & 42.3851 \\
Shuffled image & 53.9806 & 42.0954 \\
\hline
\end{tabular}
\end{table}

Shuffling the camera input reduces moderate BEV AP$_{40}$ by 0.6672 points
and moderate 3D AP$_{40}$ by 0.5742 points. It changes the predictions in
3,716 of the 3,769 validation frames, demonstrating that image information
influences the detector. The strongest class-specific effect is observed for
Pedestrian, whose moderate 3D AP$_{40}$ decreases from 32.2419 to 30.4090
under image shuffling. The corresponding changes for Car and Cyclist are
small and are not consistently positive. Thus, useful camera information is
most evident for pedestrian detection, while LiDAR remains the dominant
modality overall.

\subsection{Attention Visualization}

Figure~\ref{fig:attention_maps} visualizes a head-averaged, last-layer
attention rollout from a selected LiDAR detection token through the four
shared bottleneck tokens to the geometrically aligned camera tokens. One
moderate-difficulty true positive is shown for each class. Examples are
selected using a fixed visualization-quality criterion that combines
detection confidence, object visibility, and attention overlap in both
camera and BEV space.
Camera-invisible BEV cells are excluded, and each map is normalized
independently for visualization; therefore, color magnitude should not be
compared across examples.

\begin{figure}[ht]
\centering
\includegraphics[width=\linewidth]
{evaluation_assets/attention_maps_camera_bev.png}
\caption{Attention rollout for Pedestrian, Cyclist, and Car examples. The
left column projects attention onto the camera image, while the right column
shows the same attention in BEV together with LiDAR points and detections.
White boxes identify the selected ground-truth object and corresponding
prediction.}
\label{fig:attention_maps}
\end{figure}

The maps show nonuniform but generally diffuse attention over camera-visible
scene regions. The top 10\% of valid camera cells contain 21.04\%, 19.39\%,
and 37.57\% of the attention mass for the displayed Pedestrian, Cyclist, and
Car examples, respectively. Attention peaks overlap the selected objects in
these examples while also covering nearby objects and scene context. This
behavior is expected from four global bottleneck tokens, which compress
scene-level information rather than enforcing object-localized
correspondence. Because the examples were selected partly for attention
overlap, they are illustrative rather than representative. The maps visualize
active attention pathways but do not by themselves demonstrate causal camera
use; the camera-input ablation provides the stronger evidence.

\subsection{Discussion and Limitations}

The results show that the proposed attention-bottleneck fusion detector
performs better than the LiDAR-only reference on the KITTI validation split.
Camera ablations demonstrate a measurable contribution from visual content,
particularly for pedestrians. However, the ablation gap is smaller than the
total improvement over LiDAR-only, so the full gain cannot be attributed to
camera information alone. The additional attention layers, bottleneck
representation, LiDAR-side processing, warm start, and staged fine-tuning may
also contribute.

The experiments use one random seed and the validation split rather than the
official KITTI test server. Consequently, variation across training runs has
not been measured. Additional seeds and comparisons with alternative fusion
architectures are required before claiming that attention bottlenecks are
generally superior to other fusion methods. Attention weights should also be
interpreted as model diagnostics rather than causal explanations.
